% GENERATED FILE. Edit canonical lessons, then run npm run generate:books.
% canonical-source-hash: fnv1a64:7f7bdffee173dcb2
% canonical-lessons: HI-C388-adhik, HI-C388-tadad, HI-C388-paryapt, HI-C388-mamuli, HI-C388-qarib

\chapter{More and Less}
\label{ch:hi-more-and-less}
\hlchaptermodality{388}

\begin{chapteropening}
\textbf{By the end of this chapter:} \emph{I can say how much will be enough if things go to plan.}

Complete the last lesson of chapter 388: I can say how much will be enough if things go to plan.
\end{chapteropening}

\section[\texorpdfstring{\hi{अधिक} (adhik)}{अधिक (adhik)}]{\hi{अधिक} (adhik) — more (formal), excessive}

\label{lesson:HI-C388-adhik}



\begin{quote}
Before the new one: say the Hindi for plenty, then the Hindi for tremendous.
\end{quote}

\subsection*{You'll want to know: \hi{अधिक}}
\textbf{\hi{अधिक}} — \emph{adhik} — \textquotedblleft{}more (formal), excessive\textquotedblright{}.

\begin{quote}
\hi{अगर} \hi{पानी} \hi{अधिक} \hi{होगा} \hi{तो} \hi{चावल} \hi{गीला} \hi{हो} \hi{जाएगा।}
\end{quote}

\begin{quote}
\emph{agar pānī adhik hogā to cāval gīlā ho jāegā.}
\end{quote}

\begin{quote}
If there is too much water, the rice will turn soggy.
\end{quote}

\begin{cousinweb}
\textbf{\hi{अधिक}} (\emph{adhik}) is a tatsama, borrowed from Sanskrit \emph{adhika}, \textquotedblleft{}surpassing, more\textquotedblright{} (Monier-Williams).
\end{cousinweb}

\begin{grammarlens}[title={What you've built}]
One more word for saying what happens if.
\end{grammarlens}

\subsection*{Your turn}
\begin{itemize}
\raggedright
  \item \emph{Say it:} \emph{adhik}
  \item \emph{Say it:} \emph{adhik}, once more
  \item \emph{Say it:} \emph{agar pānī adhik hogā to cāval gīlā ho jāegā.} — If there is too much water, the rice will turn soggy.
  \item \emph{Recall:} say the Hindi for plenty, then the Hindi for tremendous, then say \emph{adhik} again
\end{itemize}

\begin{tcolorbox}[breakable,colback=teal!4,colframe=teal!35!black,title={Before you move on}]
What does \hi{अधिक} mean? (\textquotedblleft{}More (formal), excessive\textquotedblright{}.) Say it once more.
\end{tcolorbox}
\section[\texorpdfstring{\hi{तादाद} (tādād)}{तादाद (tādād)}]{\hi{तादाद} (tādād) — a number, a quantity (of people or things)}

\label{lesson:HI-C388-tadad}



\begin{quote}
Before the new one: say the Hindi for tremendous, then the Hindi for more (formal).
\end{quote}

\subsection*{You'll want to know: \hi{तादाद}}
\textbf{\hi{तादाद}} — \emph{tādād} — \textquotedblleft{}a number, a quantity (of people or things)\textquotedblright{}.

\begin{quote}
\hi{अगर} \hi{मेहमानों} \hi{की} \hi{तादाद} \hi{बढ़ेगी} \hi{तो} \hi{खाना} \hi{कम} \hi{पड़ेगा।}
\end{quote}

\begin{quote}
\emph{agar mehmānoṁ kī tādād baṛhegī to khānā kam paṛegā.}
\end{quote}

\begin{quote}
If the number of guests goes up, the food will run short.
\end{quote}

Feminine.

\begin{cousinweb}
\textbf{\hi{तादाद}} (\emph{tādād}) is usually traced to Arabic \emph{taʿdād}, \textquotedblleft{}counting, a number\textquotedblright{}, which reached Hindi through Persian.
\end{cousinweb}

\begin{grammarlens}[title={What you've built}]
One more word for saying what happens if.
\end{grammarlens}

\subsection*{Your turn}
\begin{itemize}
\raggedright
  \item \emph{Say it:} \emph{tādād}
  \item \emph{Say it:} \emph{tādād}, once more
  \item \emph{Say it:} \emph{agar mehmānoṁ kī tādād baṛhegī to khānā kam paṛegā.} — If the number of guests goes up, the food will run short.
  \item \emph{Recall:} say the Hindi for tremendous, then the Hindi for more (formal), then say \emph{tādād} again
\end{itemize}

\begin{tcolorbox}[breakable,colback=teal!4,colframe=teal!35!black,title={Before you move on}]
What does \hi{तादाद} mean? (\textquotedblleft{}A number, a quantity (of people or things)\textquotedblright{}.) Say it once more.
\end{tcolorbox}
\section[\texorpdfstring{\hi{पर्याप्त} (paryāpt)}{पर्याप्त (paryāpt)}]{\hi{पर्याप्त} (paryāpt) — sufficient}

\label{lesson:HI-C388-paryapt}



\begin{quote}
Before the new one: say the Hindi for more (formal), then the Hindi for a quantity.
\end{quote}

\subsection*{You'll want to know: \hi{पर्याप्त}}
\textbf{\hi{पर्याप्त}} — \emph{paryāpt} — \textquotedblleft{}sufficient\textquotedblright{}.

\begin{quote}
\hi{अगर} \hi{खाना} \hi{पर्याप्त} \hi{नहीं} \hi{होगा} \hi{तो} \hi{और} \hi{बनाना} \hi{पड़ेगा।}
\end{quote}

\begin{quote}
\emph{agar khānā paryāpt nahīṃ hogā to aur banānā paṛegā.}
\end{quote}

\begin{quote}
If there isn't enough food, we'll have to make more.
\end{quote}

\begin{cousinweb}
\textbf{\hi{पर्याप्त}} (\emph{paryāpt}) is a tatsama, borrowed from Sanskrit \emph{paryāpta}, \textquotedblleft{}complete; sufficient\textquotedblright{} (Monier-Williams).
\end{cousinweb}

\begin{grammarlens}[title={What you've built}]
One more word for saying what happens if.
\end{grammarlens}

\subsection*{Your turn}
\begin{itemize}
\raggedright
  \item \emph{Say it:} \emph{paryāpt}
  \item \emph{Say it:} \emph{paryāpt}, once more
  \item \emph{Say it:} \emph{agar khānā paryāpt nahīṃ hogā to aur banānā paṛegā.} — If there isn't enough food, we'll have to make more.
  \item \emph{Recall:} say the Hindi for more (formal), then the Hindi for a quantity, then say \emph{paryāpt} again
\end{itemize}

\begin{tcolorbox}[breakable,colback=teal!4,colframe=teal!35!black,title={Before you move on}]
What does \hi{पर्याप्त} mean? (\textquotedblleft{}Sufficient\textquotedblright{}.) Say it once more.
\end{tcolorbox}
\section[\texorpdfstring{\hi{मामूली} (māmūlī)}{मामूली (māmūlī)}]{\hi{मामूली} (māmūlī) — minor, ordinary}

\label{lesson:HI-C388-mamuli}



\begin{quote}
Before the new one: say the Hindi for a quantity, then the Hindi for sufficient.
\end{quote}

\subsection*{You'll want to know: \hi{मामूली}}
\textbf{\hi{मामूली}} — \emph{māmūlī} — \textquotedblleft{}minor, ordinary\textquotedblright{}.

\begin{quote}
\hi{अगर} \hi{चोट} \hi{मामूली} \hi{है} \hi{तो} \hi{पट्टी} \hi{काफ़ी} \hi{है।}
\end{quote}

\begin{quote}
\emph{agar coṭ māmūlī hai to paṭṭī kāfī hai.}
\end{quote}

\begin{quote}
If the injury is minor, a bandage is enough.
\end{quote}

\begin{cousinweb}
\textbf{\hi{मामूली}} (\emph{māmūlī}) is usually traced to Arabic \emph{maʿmūlī}, \textquotedblleft{}customary\textquotedblright{}, which reached Hindi through Persian.
\end{cousinweb}

\begin{grammarlens}[title={What you've built}]
One more word for saying what happens if.
\end{grammarlens}

\subsection*{Your turn}
\begin{itemize}
\raggedright
  \item \emph{Say it:} \emph{māmūlī}
  \item \emph{Say it:} \emph{māmūlī}, once more
  \item \emph{Say it:} \emph{agar coṭ māmūlī hai to paṭṭī kāfī hai.} — If the injury is minor, a bandage is enough.
  \item \emph{Recall:} say the Hindi for a quantity, then the Hindi for sufficient, then say \emph{māmūlī} again
\end{itemize}

\begin{tcolorbox}[breakable,colback=teal!4,colframe=teal!35!black,title={Before you move on}]
What does \hi{मामूली} mean? (\textquotedblleft{}Minor, ordinary\textquotedblright{}.) Say it once more.
\end{tcolorbox}
\section[\texorpdfstring{\hi{क़रीब} (qarīb)}{क़रीब (qarīb)}]{\hi{क़रीब} (qarīb) — close to, nearly}

\label{lesson:HI-C388-qarib}



\begin{quote}
Before the new one: say the Hindi for sufficient, then the Hindi for minor.
\end{quote}

\subsection*{You'll want to know: \hi{क़रीब}}
\textbf{\hi{क़रीब}} — \emph{qarīb} — \textquotedblleft{}close to, nearly\textquotedblright{}.

\begin{quote}
\hi{अगर} \hi{बस} \hi{समय} \hi{पर} \hi{चली} \hi{तो} \hi{हम} \hi{क़रीब} \hi{दो} \hi{घंटे} \hi{में} \hi{पहुँचेंगे।}
\end{quote}

\begin{quote}
\emph{agar bas samay par calī to ham qarīb do ghaṇṭe meṁ pahuñceṅge.}
\end{quote}

\begin{quote}
If the bus leaves on time, we'll get there in close to two hours.
\end{quote}

\begin{cousinweb}
\textbf{\hi{क़रीब}} (\emph{qarīb}) is usually traced to Arabic \emph{qarīb}, \textquotedblleft{}near\textquotedblright{}, which reached Hindi through Persian.
\end{cousinweb}

\begin{grammarlens}[title={What you've built}]
One more word for saying what happens if.
\end{grammarlens}

\subsection*{Your turn}
\begin{itemize}
\raggedright
  \item \emph{Say it:} \emph{qarīb}
  \item \emph{Say it:} \emph{qarīb}, once more
  \item \emph{Say it:} \emph{agar bas samay par calī to ham qarīb do ghaṇṭe meṁ pahuñceṅge.} — If the bus leaves on time, we'll get there in close to two hours.
  \item \emph{Recall:} say the Hindi for sufficient, then the Hindi for minor, then say \emph{qarīb} again
\end{itemize}

\begin{tcolorbox}[breakable,colback=teal!4,colframe=teal!35!black,title={Before you move on}]
What does \hi{क़रीब} mean? (\textquotedblleft{}Close to, nearly\textquotedblright{}.) Say it once more.
\end{tcolorbox}
